\documentclass[../../../include/open-logic-section]{subfiles}

\begin{document}

\olfileid{sth}{ordinals}{wo}
\olsection{الترتيبات الجيدة}

المفهوم الأساسي هو الآتي:

\begin{defn}
تكون العلاقة~$<$ \emph{ترتيبًا جيدًا} على~$A$ إذا وفقط إذا استوفت
الشرطين الآتيين:
\begin{enumerate}
	\item تكون~$<$ متصلة، أي لكل~$a, b \in A$، إما~$a < b$،
	أو~$a = b$، أو~$b < a$؛
	\item لكل مجموعة جزئية غير خالية من~$A$ !!{element} أصغري بحسب~$<$،
	أي إذا كان~$\emptyset \neq X \subseteq A$، فإن
	$(\exists m \in X)(\forall z \in X)z \nless m$.
\end{enumerate}
\end{defn}

يسهل التحقق من أن العلاقات في الأمثلة الثلاثة التي نظرنا فيها للتو
كانت بالفعل ترتيبات جيدة.

\begin{prob}
عرض~\Olref[sth][ordinals][idea]{sec} ثلاثة ترتيبات نموذجية على الأعداد
الطبيعية. تحقق من أن كلًا منها ترتيب جيد.
\end{prob}

وفيما يأتي بضع ملاحظات أولية، لكنها شديدة الأهمية، بشأن الترتيب الجيد.

\begin{prop}\ollabel{wo:strictorder}
إذا كانت~$<$ ترتيبًا جيدًا على~$A$، فلكل مجموعة جزئية غير خالية من~$A$
عنصر أصغر وحيد بحسب~$<$، وتكون~$<$ غير انعكاسية ولا تناظرية ومتعدية.
\end{prop}

\begin{proof}
إذا كانت~$X$ مجموعة جزئية غير خالية من~$A$، فلها !!{element} أصغري
بحسب~$<$، وليكن~$m$؛ أي $(\forall z \in X)z \nless m$. ولما كانت~$<$
متصلة، كان $(\forall z \in X)m \leq z$. ومن ثم يكون~$m$ أصغر
!!{element} بحسب~$<$ في~$X$.

ولإثبات عدم الانعكاس، ثبّت~$a \in A$؛ فأصغر !!{element} بحسب~$<$ في
$\{a\}$ هو~$a$، ولذلك~$a \nless a$. ولإثبات التعدي، إذا كان
$a < b < c$، فبما أن~$\{a, b, c\}$ لها !!{element} أصغر بحسب~$<$،
كان~$a < c$. وينتج عدم التناظر من عدم الانعكاس والتعدي.
\end{proof}

\begin{prop}\ollabel{propwoinduction}
إذا كانت~$<$ ترتيبًا جيدًا على~$A$، فعندئذ، لكل صيغة~$\phi(x)$:
\[
	\text{إذا كان }(\forall a \in A)((\forall b < a)\phi(b) \lif 
		\phi(a))\text{، فإن }(\forall a \in A)\phi(a).
\]
\end{prop}

\begin{proof}
سنثبت المقابل بالنقيض. افترض~$\lnot(\forall a \in A)\phi(a)$، أي إن
$X = \Setabs{x \in A}{\lnot\phi(x)} \neq \emptyset$. عندئذ يكون
لـ$X$ !!{element} أصغري بحسب~$<$، وليكن~$a$. ومن ثم
$(\forall b < a)\phi(b)$، مع أن~$\lnot \phi(a)$.
\end{proof}\noindent
ينبغي أن تذكرك هذه الخاصية الأخيرة بمبدأ الاستقراء القوي على الأعداد
الطبيعية، أي: إذا كان
$(\forall n \in \omega)((\forall m < n)\phi(m) \lif \phi(n))$، فإن
$(\forall n \in \omega)\phi(n)$. وهذه الخاصية هي ما يجعل الترتيب
الجيد مفهومًا بالغ \emph{المتانة}.
\footnote{تذكير: يجوز أن تكون لجميع الصيغ معاملات (ما لم يُصرح بخلاف ذلك).}

\end{document}
